\documentclass[11pt]{article}
\newcommand{\paso}{diccionario}
\usepackage[T1]{fontenc}
\usepackage[utf8]{inputenc}
\usepackage[spanish,es-nodecimaldot,es-noquoting,es-lcroman]{babel}
\spanishplainpercent
\usepackage{lmodern}
\usepackage[a4paper,margin=21mm,headheight=15pt]{geometry}
\usepackage{amsmath,amssymb,bm}
\usepackage{booktabs,array}
\usepackage{tikz}
\usetikzlibrary{arrows.meta,positioning,calc}
\usepackage[most]{tcolorbox}
\usepackage{microtype}
\usepackage{fancyhdr}
\usepackage{hyperref}
\hypersetup{colorlinks=true,linkcolor=blue!45!black,urlcolor=blue!45!black}
\setlength{\parindent}{0pt}
\setlength{\parskip}{5pt}
\setlength{\emergencystretch}{2em}
\pagestyle{fancy}
\fancyhf{}
\fancyhead[L]{\small BEM del cilindro: paso \paso}
\fancyhead[R]{\small Guía matemática independiente}
\fancyfoot[C]{\thepage}
\newcommand{\vct}[1]{\bm{#1}}
\newcommand{\dd}{\,\mathrm{d}}
\newcommand{\ii}{\mathrm{i}}
\newcommand{\norm}[1]{\left\lVert#1\right\rVert}
\newcommand{\VP}{\operatorname{VP}}
\newcommand{\tr}{\operatorname{tr}}
\newtcolorbox{idea}[1]{breakable,colback=blue!3,colframe=blue!50!black,
  title={#1},fonttitle=\bfseries,boxrule=0.7pt}
\newtcolorbox{alerta}[1]{breakable,colback=orange!4,colframe=orange!65!black,
  title={#1},fonttitle=\bfseries,boxrule=0.7pt}
\newtcolorbox{resultado}[1]{breakable,colback=teal!3,colframe=teal!55!black,
  title={#1},fonttitle=\bfseries,boxrule=0.7pt}
\newcommand{\inicio}[2]{%
 {\small\color{blue!55!black}\bfseries SERIE: LOS SIETE PASOS MATEMÁTICOS DEL BEM}\par
 {\LARGE\bfseries Paso \paso. #1\par}\vspace{3mm}
 \begin{idea}{¿Qué busca este paso?}#2\end{idea}}
\newcommand{\mapa}[3]{%
 \begin{center}
 \begin{tikzpicture}[node distance=4mm,
   caja/.style={draw=blue!50!black,fill=blue!3,rounded corners=2pt,
     text width=13cm,align=center,inner sep=6pt},
   >={Stealth[length=2mm]}]
 \node[caja] (a){#1};
 \node[caja,below=of a] (b){#2};
 \node[caja,below=of b] (c){#3};
 \draw[->,thick,blue!60!black] (a)--(b);
 \draw[->,thick,blue!60!black] (b)--(c);
 \end{tikzpicture}
 \end{center}}
\newcommand{\base}{%
 \smallskip{\footnotesize\color{black!65}
 Modelo de referencia: cilindro sólido, material homogéneo e isótropo,
 deformación plana, pequeñas deformaciones, sin fuerzas de volumen;
 amplitudes complejas con \(e^{+\ii\omega t}\), \(\omega>0\).}\par}
\newcommand{\fuentes}{%
 \par\medskip\begingroup\footnotesize
 \textbf{Referencia y continuación.}
 Este documento desarrolla una etapa del esquema
 \texttt{esquema\_ecuaciones\_bem.tex}. Las ecuaciones y los ejemplos se
 contrastaron con \texttt{bem\_cilindro.py}, ambos en la carpeta superior.
 La guía general y sus deducciones ampliadas se encuentran en
 \texttt{../guia\_bem/guia\_bem.pdf}. No se modifica la implementación.
 \par Para recompilar este documento, mantener \texttt{estilo.tex} en
 la misma carpeta y ejecutar \texttt{pdflatex} sobre su archivo
 \texttt{.tex}, desde esa carpeta.\par\endgroup}

\usepackage{multicol,needspace}
\fancyhead[L]{\small BEM del cilindro: diccionario de términos}
\hypersetup{pdftitle={Diccionario visual de términos del BEM},
  pdfsubject={Conceptos para interpretar los siete pasos matemáticos}}

\newcommand{\entrada}[4]{%
  \par\addvspace{8pt}\needspace{7\baselineskip}%
  \phantomsection\label{term:#1}%
  \noindent{\large\bfseries\color{blue!55!black}#2}%
  \hfill{\footnotesize\bfseries\color{teal!55!black}Pasos #3}\par\nobreak
  \noindent#4\par}
\newcommand{\indice}[2]{%
  \noindent\hyperref[term:#1]{#2}\nobreak\dotfill\pageref{term:#1}\par}
\newcommand{\grupo}[1]{%
  \par\addvspace{8pt}\needspace{8\baselineskip}%
  {\Large\bfseries\color{blue!55!black}#1}\par}

\begin{document}
{\small\color{blue!55!black}\bfseries
SERIE: LOS SIETE PASOS MATEMÁTICOS DEL BEM}\par
{\LARGE\bfseries Diccionario visual de términos del BEM\par}
\vspace{2mm}
{\large Una guía para leer las palabras, los símbolos y las conexiones
del problema del cilindro.}\par

\begin{idea}{¿Para qué sirve este diccionario?}
Cuando aparezca una palabra como \emph{fuente}, \emph{núcleo},
\emph{global} o \emph{dinámico}, busca su definición y el número de los
pasos donde se utiliza. Cada entrada explica la idea con palabras y,
cuando ayuda, la acompaña con una fórmula o un ejemplo del cilindro.
\end{idea}

\section*{El recorrido en una sola página}
\begin{center}
\begin{tikzpicture}[node distance=2.2mm,
 caja/.style={draw=blue!50!black,fill=blue!3,rounded corners=2pt,
   text width=14.5cm,align=left,inner sep=4pt,font=\small},
 >={Stealth[length=2mm]}]
 \node[caja] (a) {\textbf{1. Problema físico:}
   cilindro \(\Omega\), frontera \(\Gamma\), presión y equilibrio.};
 \node[caja,below=of a] (b) {\textbf{2. Material:}
   \(\rho,\nu,\beta\longrightarrow\mu^*,\lambda^*,c_s,c_p\).};
 \node[caja,below=of b] (c) {\textbf{3. Identidad de Green:}
   relacionar desplazamiento y tracción en la frontera.};
 \node[caja,below=of c] (d) {\textbf{4. Núcleos:}
   \(U,T\) describen la influencia entre fuente e integración.};
 \node[caja,below=of d] (e) {\textbf{5. Discretización:}
   elementos, nodos, interpolación y cuadratura.};
 \node[caja,below=of e] (f) {\textbf{6. Ensamblaje:}
   sumar influencias y construir \(H(\omega)\) y \(G(\omega)\).};
 \node[caja,below=of f] (g) {\textbf{7. Solución:}
   aplicar \(\vct t=-p\vct n\), resolver \(H\vct u=G\vct t\) y obtener \(u_r\).};
 \foreach \x/\y in {a/b,b/c,c/d,d/e,e/f,f/g}
   \draw[->,thick,blue!60!black] (\x)--(\y);
\end{tikzpicture}
\end{center}

\begin{resultado}{Convención común}
El ejemplo es un cilindro \textbf{sólido}, homogéneo e isótropo, en
deformación plana, sin fuerzas de volumen, sometido a presión armónica.
Las amplitudes complejas usan \(e^{+\ii\omega t}\), con \(\omega>0\).
Los números de paso remiten a los siete documentos de la misma carpeta.
\end{resultado}

\newpage
\section*{Índice alfabético}
Pulsa un término para ir a su definición. Los nombres alternativos remiten
a la misma entrada; el número indica la página de este documento.
\begin{multicols}{2}
\small\setlength{\parskip}{0pt}
\indice{amortiguamiento}{Amortiguamiento}
\indice{amplitud}{Amplitud}
\indice{armonico}{Armónico}
\indice{bem}{BEM}
\indice{bessel}{Bessel}
\indice{bloque}{Bloque}
\indice{cambiovariable}{Cambio de variable}
\indice{campo}{Campo}
\indice{carga}{Carga}
\indice{poisson}{Coeficiente de Poisson}
\indice{colocacion}{Colocación}
\indice{condicion}{Condición de contorno}
\indice{conectividad}{Conectividad}
\indice{contorno}{Contorno}
\indice{coordenada}{Coordenada local}
\indice{correccion}{Corrección dinámica}
\indice{cuadratura}{Cuadratura}
\indice{deformacion}{Deformación}
\indice{plana}{Deformación plana}
\indice{delta}{Delta de Dirac y de Kronecker}
\indice{densidad}{Densidad}
\indice{derivadas}{Derivada, gradiente y divergencia}
\indice{desplazamiento}{Desplazamiento}
\indice{radial}{Desplazamiento radial}
\indice{dinamico}{Dinámico}
\indice{direccion}{Dirección fuente--punto}
\indice{discretizacion}{Discretización}
\indice{distancia}{Distancia entre puntos}
\indice{dominio}{Dominio}
\indice{ecuacioncampo}{Ecuación de campo}
\indice{ecuacionintegral}{Ecuación integral}
\indice{elemento}{Elemento de contorno}
\indice{ensamblaje}{Ensamblaje}
\indice{esfuerzo}{Esfuerzo}
\indice{estatico}{Estático}
\indice{fase}{Fase}
\indice{frecuencia}{Frecuencia}
\indice{contorno}{Frontera}
\indice{fuente}{Fuente}
\indice{fuerzaauxiliar}{Fuerza puntual auxiliar}
\indice{forma}{Función de forma}
\indice{funcionesradiales}{Funciones radiales}
\indice{isoparametrico}{Geometría isoparamétrica}
\indice{global}{Global}
\indice{green}{Identidad de Green}
\indice{incognita}{Incógnita y dato conocido}
\indice{indice}{Índice}
\indice{inercia}{Inercia}
\indice{integracion}{Integración}
\indice{interpolacion}{Interpolación}
\indice{jacobiano}{Jacobiano}
\indice{kelvin}{Kelvin}
\indice{lame}{Lamé}
\indice{local}{Local}
\indice{malla}{Malla}
\indice{material}{Material complejo}
\indice{matrices}{Matrices \(H\) y \(G\)}
\indice{medioinfinito}{Medio infinito auxiliar}
\indice{rigido}{Movimiento rígido}
\indice{nodo}{Nodo}
\indice{normal}{Normal exterior}
\indice{nucleo}{Núcleo \(U\) y \(T\)}
\indice{numeroonda}{Número de onda}
\indice{ondas}{Ondas P y S}
\indice{operador}{Operador}
\indice{particion}{Partición de la unidad}
\indice{peso}{Peso de cuadratura}
\indice{presion}{Presión}
\indice{punto}{Punto de integración}
\indice{reciprocidad}{Reciprocidad}
\indice{regularizacion}{Regularización}
\indice{residuo}{Residuo algebraico}
\indice{resonancia}{Resonancia}
\indice{singularidad}{Singularidad}
\indice{sistema}{Sistema lineal}
\indice{solucionfundamental}{Solución fundamental}
\indice{tangente}{Tangente}
\indice{tensor}{Tensor}
\indice{terminolibre}{Término libre}
\indice{traccion}{Tracción}
\indice{valorprincipal}{Valor principal}
\indice{vector}{Vector}
\indice{velocidad}{Velocidad de onda}
\end{multicols}

\newpage
\grupo{A--C}

\entrada{amortiguamiento}{Amortiguamiento}{2, 7}{
Es el mecanismo que representa la disipación de energía y produce
desfase. En este modelo se introduce en el módulo de corte como
\(\mu^*=\mu_0(1+2\ii\beta)\). Con \(\beta=0.05\), la parte
imaginaria vale el \(10\%\) de \(\mu_0\). No es lo mismo que la
densidad ni que la frecuencia.}

\entrada{amplitud}{Amplitud}{1, 7}{
Es el tamaño de una oscilación. Si una respuesta radial compleja es
\(u_r=a+\ii b\), su amplitud es \(|u_r|=\sqrt{a^2+b^2}\), siempre no
negativa. La amplitud por sí sola no indica si en \(t=0\) el punto se
mueve hacia dentro o hacia fuera: para eso también se necesita la fase.}

\entrada{armonico}{Armónico}{1, 2, 7}{
Significa que se estudia una variación sinusoidal a una frecuencia fija.
Se representa mediante una amplitud compleja \(u(\omega)\) y la regla
\(u_{\text{fís}}(t)=\operatorname{Re}\{u(\omega)e^{\ii\omega t}\}\).
Se repite el cálculo si se desea conocer la respuesta a otras frecuencias.}

\entrada{bem}{BEM: método de elementos de contorno}{3--6}{
Método numérico que formula el problema mediante integrales sobre el
contorno y discretiza esa frontera en elementos. Para este cilindro
bidimensional se discretiza la circunferencia, no el área interior.
Se obtiene un sistema denso; la reducción de dimensión viene acompañada
de núcleos e integrales singulares.}

\entrada{bessel}{Bessel: funciones \(J_n\) y \(K_n\)}{4, 7}{
Son funciones especiales que aparecen al resolver ecuaciones de onda.
En los núcleos se usan \(K_n\) con argumentos complejos; en la solución
analítica radial se usan \(J_0,J_1\). El símbolo \(J\) sin subíndice
también se usa para el jacobiano geométrico: son objetos distintos.}

\entrada{bloque}{Bloque de matriz}{6}{
Es una pequeña submatriz que reúne varias componentes relacionadas.
Aquí \(H_{cg}\) y \(G_{cg}\) son bloques \(2\times2\): las filas
corresponden a las dos ecuaciones de la fuente \(c\), y las columnas,
a las dos componentes del nodo global \(g\). No es una sola entrada
escalar.}

\entrada{cambiovariable}{Cambio de variable}{5}{
Es reescribir una integral mediante una nueva coordenada y su factor de
escala. Cerca de una fuente situada en el extremo \(a\), se puede usar
\(\eta=a+(b-a)s^2\), con \(0<s<1\). Esto concentra puntos cerca de la
fuente; el peso de integración también debe incluir
\(\dd\eta/\dd s=2(b-a)s\).}

\entrada{campo}{Campo}{1, 3}{
Es una cantidad definida en muchos puntos del espacio. Por ejemplo,
\(\vct u(\vct x)\) asigna un desplazamiento a cada punto \(\vct x\);
\(\boldsymbol\sigma(\vct x)\) asigna un tensor de esfuerzo. En BEM
interesan especialmente sus valores sobre el contorno.}

\entrada{carga}{Carga}{1, 7}{
Es la acción aplicada al cuerpo. En el ejemplo es una presión uniforme
sobre el contorno, representada por la tracción \(\vct t=-p\vct n\).
Una \emph{fuerza de volumen}, como el peso por unidad de volumen,
actuaría en el interior; este modelo no la incluye.}

\entrada{poisson}{Coeficiente de Poisson \(\nu\)}{1, 2}{
Relaciona deformaciones laterales y longitudinales en elasticidad
isótropa. En el ejemplo \(\nu=0.25\) y se usa
\(\lambda^*=2\nu\mu^*/(1-2\nu)\). No representa amortiguamiento:
ese papel corresponde a \(\beta\).}

\entrada{colocacion}{Colocación}{3, 6}{
Consiste en elegir puntos donde se impone la ecuación integral.
Aquí se usa cada nodo global del contorno como fuente
\(\vct\xi_c=\vct x_c\). En cada uno se escriben dos ecuaciones, una
por componente cartesiana. Colocación no significa colocar físicamente
la presión en un solo nodo.}

\entrada{condicion}{Condición de contorno}{1, 7}{
Es el dato prescrito en la frontera. Puede ser desplazamiento, tracción
u otra condición, según el problema. En este caso se conoce la tracción
en todo el borde y se busca el desplazamiento allí. No se prescriben
ambos campos completos como datos independientes en el mismo borde.}

\entrada{conectividad}{Conectividad}{5, 6}{
Es la lista que dice qué nodos globales forman cada elemento.
Con índices globales desde cero, los dos primeros elementos usan
\((0,1,2)\) y \((2,3,4)\). El nodo 2 es compartido: sus contribuciones
de ambos elementos se suman en el ensamblaje.}

\entrada{contorno}{Contorno o frontera \(\Gamma\)}{1, 3, 5}{
Es la curva que limita la región material. En este problema es la
circunferencia de radio \(R\). La ecuación BEM integra sobre esa curva.
El contorno es el lugar geométrico; la condición de contorno es el dato
que se impone allí.}

\entrada{coordenada}{Coordenada local \(\eta\)}{5}{
Es una posición adimensional dentro de un elemento de referencia:
\(-1\le\eta\le1\). Los tres nodos locales están en \(-1,0,1\).
La geometría interpolada transforma \(\eta\) en un punto físico
\(\vct x(\eta)\); \(\eta\) no es una distancia en metros.}

\entrada{correccion}{Corrección dinámica \(\Delta H\)}{4, 6}{
Es el aporte que se suma a la matriz estática:
\(H(\omega)=H^s+\Delta H(\omega)\). Se integra usando \(T-T^s\), incluida
la diagonal. No debe imponerse a \(H(\omega)\) la suma de bloques cero
que se usa para construir \(H^s\).}

\entrada{cuadratura}{Cuadratura}{5, 6}{
Es una regla para aproximar una integral mediante una suma de valores
en puntos y pesos seleccionados. Aquí se usa Gauss-Legendre; cuando la
fuente pertenece al elemento se transforma la variable y se concentran
puntos cerca de ella. Evitar la fuente no basta, por sí solo, para
resolver un valor principal.}

\grupo{D--F}

\entrada{deformacion}{Deformación \(\boldsymbol\varepsilon\)}{1}{
Mide el cambio local de forma o tamaño, no el desplazamiento total.
Para pequeñas deformaciones,
\(\varepsilon_{ij}=\tfrac12(u_{i,j}+u_{j,i})\). Una traslación
constante tiene desplazamiento distinto de cero, pero deformación cero.}

\entrada{plana}{Deformación plana}{1}{
Hipótesis bidimensional con \(u_z=0\), sin dependencia de \(z\), y
deformaciones axiales y cortantes con \(z\) nulas. Puede existir
\(\sigma_{zz}\). No equivale a \emph{esfuerzo plano}, donde se impone
\(\sigma_{zz}=0\).}

\entrada{delta}{Delta de Dirac y delta de Kronecker}{3, 4}{
\(\delta(\vct x-\vct\xi)\) representa una fuente puntual en una integral;
la integral de un campo continuo contra ella selecciona su valor en
\(\vct\xi\). En cambio, \(\delta_{ij}\) vale 1 si \(i=j\) y 0 si no.
Comparten letra, pero no son el mismo objeto.}

\entrada{densidad}{Densidad \(\rho\)}{1, 2}{
Es la masa por unidad de volumen, en \(\mathrm{kg/m^3}\). Multiplica
la aceleración en el equilibrio dinámico y participa en las velocidades
de onda. El ejemplo usa \(\rho=100\ \mathrm{kg/m^3}\).}

\entrada{derivadas}{Derivada, gradiente y divergencia}{1, 3, 4}{
Una derivada mide cómo cambia una cantidad al cambiar la posición.
El gradiente reúne derivadas espaciales; la divergencia contrae una
derivada con un índice del campo. Por ejemplo,
\((\nabla\!\cdot\boldsymbol\sigma)_i=\sigma_{ij,j}\) da la fuerza
interna neta por unidad de volumen.}

\entrada{desplazamiento}{Desplazamiento \(\vct u\)}{1, 7}{
Es el cambio de posición respecto a la configuración de referencia.
En el plano tiene componentes \((u_x,u_y)\), medidas en metros.
La letra minúscula \(\vct u\) es el campo real buscado en amplitudes
complejas; la mayúscula \(U\) es un núcleo auxiliar.}

\entrada{radial}{Desplazamiento radial y tangencial}{7}{
Son las componentes del desplazamiento en la dirección del radio y en
la dirección perpendicular. En el nodo circular,
\(u_r=\vct u\cdot\vct n\). Para presión uniforme y simetría perfecta,
la solución exacta tiene componente tangencial nula.}

\entrada{dinamico}{Dinámico}{1, 2, 4, 6, 7}{
Describe una respuesta que depende del tiempo y tiene en cuenta la
inercia. Aquí se estudia una frecuencia \(\omega\) por vez, mediante
amplitudes complejas. Los núcleos \(U,T\) y las matrices \(H,G\)
dependen de \(\omega\). Contrasta con el estado estático auxiliar.}

\entrada{direccion}{Dirección fuente--punto \(\vct e\)}{4}{
Es el vector unitario que apunta desde la fuente hasta el punto de
integración: \(\vct e=(\vct x-\vct\xi)/r\).
Depende de ambos puntos y aparece en los núcleos.
No debe confundirse con la normal exterior \(\vct n\) ni con la dirección
radial del cilindro, que se mide desde su centro.}

\entrada{discretizacion}{Discretización}{5, 6}{
Es sustituir una descripción continua por un número finito de piezas y
valores. Se aproxima la frontera con elementos y los campos mediante
valores nodales. La solución discreta puede tener error aunque el sistema
lineal se resuelva con un residuo muy pequeño.}

\entrada{distancia}{Distancia fuente--integración \(r\)}{4}{
Para una fuente fija \(\vct\xi\) y un punto \(\vct x\),
\(r=\norm{\vct x-\vct\xi}\). Cambia al recorrer el contorno y puede
acercarse a cero. No es el radio fijo \(R\) del cilindro.}

\entrada{dominio}{Dominio \(\Omega\)}{1, 3}{
Es la región ocupada por el material que se estudia. Aquí es el disco
\(x^2+y^2<R^2\); su borde es \(\Gamma\). Un dominio no es lo mismo que
el medio infinito auxiliar usado para construir las soluciones
fundamentales.}

\entrada{ecuacioncampo}{Ecuación de campo o equilibrio}{1}{
Expresa lo que debe cumplirse en el interior. En amplitudes armónicas
sin fuerzas de volumen:
\(\nabla\!\cdot\boldsymbol\sigma+\rho\omega^2\vct u=\vct0\).
Combina esfuerzos e inercia; junto con las condiciones de contorno
define el problema físico.}

\entrada{ecuacionintegral}{Ecuación integral de contorno}{3, 6}{
Relaciona los valores de desplazamiento y tracción sobre la frontera:
\(c\vct u+\VP\int_\Gamma T\vct u\,\dd\Gamma
=\int_\Gamma U\vct t\,\dd\Gamma\).
El término con \(T\) contiene la incógnita del borde; el de \(U\)
contiene la tracción conocida en este ejemplo.}

\entrada{elemento}{Elemento de contorno}{5}{
Es un tramo de la frontera discretizada. Aquí cada elemento es una
curva interpolada con tres nodos locales: dos extremos y uno intermedio.
Los nodos extremos se comparten con elementos vecinos; por eso
ocho elementos cerrados producen 16 nodos globales, no 24.}

\entrada{ensamblaje}{Ensamblaje}{6}{
Es reunir en matrices globales las contribuciones calculadas para cada
fuente, elemento y nodo local. Si dos elementos comparten un nodo, sus
aportes se \textbf{suman} en la misma columna global. Ensamblar no es
resolver el sistema: eso ocurre después.}

\entrada{esfuerzo}{Esfuerzo \(\boldsymbol\sigma\)}{1, 3}{
Describe fuerzas internas por unidad de área y se expresa como tensor
en Pa. La tracción sobre una superficie concreta se obtiene al proyectar
ese tensor con su normal: \(\vct t=\boldsymbol\sigma\vct n\).
Por tanto, esfuerzo y tracción no son sinónimos.}

\entrada{estatico}{Estático}{1, 4, 6}{
Describe un estado sin aceleración. En el procedimiento se usa el
núcleo estático de Kelvin \(T^s\) como referencia para tratar la
singularidad, aunque el problema final sea dinámico. No significa que
se sustituya directamente \(\omega=0\) en los núcleos dinámicos del código.}

\entrada{fase}{Fase \(\phi\)}{1, 2, 7}{
Indica el desplazamiento temporal de una oscilación respecto de la
referencia elegida. Si \(u_r=Ae^{\ii\phi}\), la respuesta física es
\(A\cos(\omega t+\phi)\). Para comparar fases hay que mantener la misma
convención temporal y el mismo signo de referencia de la carga.}

\entrada{frecuencia}{Frecuencia \(f\) y frecuencia angular \(\omega\)}{1, 2, 4, 7}{
\(f\) cuenta ciclos por segundo y se mide en Hz. La frecuencia angular
se mide en rad/s y satisface \(\omega=2\pi f\). Los núcleos y el sistema
dinámico se calculan para cada \(\omega\); en el ejemplo se usa
\(\omega=40\ \mathrm{rad/s}\).}

\entrada{fuente}{Fuente \(\vct\xi\)}{3, 4, 6}{
Es el punto fijo que identifica la fuerza puntual del \emph{problema
auxiliar} y la posición para la cual se escribe la ecuación integral.
En colocación coincide con un nodo del contorno. No es la presión real
aplicada al cilindro; el punto de integración \(\vct x\) recorre la
frontera mientras la fuente permanece fija.}

\entrada{fuerzaauxiliar}{Fuerza puntual auxiliar}{3, 4}{
Es una fuerza unitaria idealizada aplicada en la fuente para definir
una solución fundamental del medio. Su función es matemática: construir
influencias. La carga física del ejemplo sigue siendo la presión
distribuida sobre todo el contorno.}

\entrada{forma}{Función de forma \(P_a\)}{5, 6}{
Es una función que asigna peso a un nodo local al interpolar geometría,
desplazamiento o tracción. En el elemento cuadrático:
\(P_1=\eta(\eta-1)/2\), \(P_2=1-\eta^2\),
\(P_3=\eta(\eta+1)/2\). En cada nodo local, su propia función vale 1
y las otras dos valen 0.}

\entrada{funcionesradiales}{Funciones radiales \(\psi,\chi\)}{4}{
Son funciones escalares de la distancia y la frecuencia que agrupan las
funciones de Bessel en el núcleo de desplazamiento:
\(U_{ij}=(\psi\delta_{ij}-\chi e_ie_j)/(2\pi\mu^*)\).
Se calculan con \(z_s=\ii\omega r/c_s\), \(z_p=\ii\omega r/c_p\) y
\(\gamma=c_s/c_p\). No son funciones de forma del elemento.}

\grupo{G--N}

\entrada{isoparametrico}{Geometría isoparamétrica}{5}{
Significa que las mismas funciones de forma se usan para interpolar la
geometría y los campos. Los nodos están sobre el círculo, pero la curva
polinómica entre ellos no es un arco circular exacto. Refinar la malla
reduce esta aproximación geométrica.}

\entrada{global}{Global}{5, 6}{
Se refiere a la numeración o al sistema de \textbf{toda} la malla, no
a un solo elemento. Por ejemplo, el nodo global \(g=2\) es el tercer
nodo local del elemento 0 y el primer nodo local del elemento 1.
Las matrices globales reúnen los aportes de todos los elementos.}

\entrada{green}{Identidad de Green o de contorno}{3}{
Es la relación matemática que transforma el equilibrio en una expresión
de integrales sobre la frontera usando un estado auxiliar. Su versión
interior se conoce como representación de Somigliana. Permite reducir
la dimensión de la discretización, no elimina la física interior.}

\entrada{incognita}{Incógnita y dato conocido}{1, 6, 7}{
Un dato conocido se prescribe; una incógnita debe calcularse. Aquí
\(\vct t=-p\vct n\) es conocida y el desplazamiento nodal \(\vct u\)
es desconocido. Una vez elegida la frecuencia, \(H\) y \(G\) se
calculan antes de resolver \(H\vct u=G\vct t\).}

\entrada{indice}{Índice}{4--6}{
Es una etiqueta que indica qué objeto o componente se usa.
\(c\): nodo de colocación; \(g\): nodo global; \(e\): elemento;
\(a\): nodo local; \(q\): punto de cuadratura;
\(i,j\): componentes cartesianas. La misma letra puede tener otro
significado en otro contexto; mira siempre su definición.}

\entrada{inercia}{Inercia}{1, 2, 7}{
Es la oposición de la masa a acelerar. La ecuación temporal contiene
\(\rho\ddot{\vct u}\); en amplitudes con \(e^{+\ii\omega t}\),
aparece el término \(+\rho\omega^2\vct u\) en el lado izquierdo del
equilibrio. No es una fuerza de amortiguamiento.}

\entrada{integracion}{Integración sobre el contorno}{3, 5, 6}{
Es acumular contribuciones de todos los puntos de \(\Gamma\). En BEM
cada fuente recibe la influencia de todo el borde. Tras discretizar,
la integral se separa por elementos y se aproxima mediante cuadratura.
La integral continua y la suma numérica no son idénticas en general.}

\entrada{interpolacion}{Interpolación}{5, 6}{
Es aproximar un valor entre nodos a partir de sus valores nodales.
Por ejemplo,
\(\vct u_e(\eta)\simeq\sum_{a=1}^3P_a(\eta)\vct u_{g(e,a)}\).
La interpolación describe cómo varía el campo; la cuadratura aproxima
la integral de ese campo.}

\entrada{jacobiano}{Jacobiano geométrico \(J\)}{5}{
Es el factor que convierte un incremento de la coordenada local en
longitud sobre el elemento:
\(J=\norm{\dd\vct x/\dd\eta}\) y
\(\dd\Gamma=J\dd\eta\). Tiene unidades de metros porque \(\eta\) es
adimensional. No confundirlo con las funciones de Bessel \(J_0,J_1\).}

\entrada{kelvin}{Kelvin: núcleo estático \(T^s\)}{4, 6}{
Es la solución fundamental estática de un medio elástico homogéneo.
Su tracción \(T^s\) reproduce la parte singular fuerte del núcleo
dinámico \(T\). Se usa en la separación \(T=T^s+(T-T^s)\) y en la
construcción de \(H^s\).}

\entrada{lame}{Lamé: parámetros \(\lambda^*,\mu^*\)}{1, 2}{
Son módulos que describen la rigidez de un material elástico isótropo.
\(\mu^*\) es el módulo de corte y \(\lambda^*\) interviene en la parte
volumétrica de la ley de Hooke; ambos se miden en Pa. El asterisco
indica que pueden ser complejos en el modelo amortiguado.}

\entrada{local}{Local}{5, 6}{
Se refiere a un solo elemento o a su coordenada y nodos internos.
Cada elemento tiene nodos locales \(a=1,2,3\) y coordenada
\(\eta\in[-1,1]\). La conectividad traduce cada par \((e,a)\) a un
nodo global \(g(e,a)\). En el código los índices locales empiezan en 0.}

\entrada{malla}{Malla de contorno}{5}{
Es el conjunto organizado de nodos y elementos que aproxima
\(\Gamma\). Refinar la malla suele significar usar más elementos y
nodos. Es diferente de aumentar los puntos de cuadratura dentro de
cada elemento.}

\entrada{material}{Material complejo}{2}{
Es la descripción armónica en la que los módulos tienen parte real e
imaginaria. La parte imaginaria de \(\mu^*\) representa el desfase
disipativo bajo la convención adoptada. No indica que el material físico
ni su desplazamiento medible sean números complejos.}

\entrada{matrices}{Matrices \(H\) y \(G\)}{6, 7}{
Son tablas de coeficientes que reúnen las influencias de todos los nodos.
\(H\) multiplica los desplazamientos desconocidos y \(G\) las tracciones
conocidas: \(H\vct u=G\vct t\). Con \(N\) nodos ambas tienen tamaño
\(2N\times2N\); son densas y, en general, complejas.}

\entrada{medioinfinito}{Medio infinito auxiliar}{3, 4}{
Es el medio ideal sin frontera en el que se define la respuesta a la
fuerza puntual auxiliar. Tiene las propiedades del material del problema.
No convierte el cilindro en un sólido infinito: la ecuación integral
impone después las condiciones del contorno real.}

\entrada{rigido}{Movimiento o traslación rígida}{6}{
Una traslación rígida desplaza todos los puntos por el mismo vector
sin producir deformación ni esfuerzo. En el problema \emph{estático}
auxiliar permite obtener
\(H^s_{cc}=-\sum_{g\ne c}H^s_{cg}\). Esta suma es de bloques
\(2\times2\), no de todas las entradas escalares.}

\entrada{nodo}{Nodo}{5, 6}{
Es un punto donde se guardan valores discretos del campo. Cada nodo
global tiene dos componentes de desplazamiento. Un nodo de colocación
se usa además como fuente; un punto de cuadratura se usa para integrar
y normalmente no es un nodo.}

\entrada{normal}{Normal exterior \(\vct n\)}{1, 5, 7}{
Es un vector unitario perpendicular al borde y dirigido fuera del
material. La presión positiva actúa hacia dentro: \(\vct t=-p\vct n\).
Para integrar se usa la normal de la curva interpolada; en los nodos,
para presión y proyección, se usa \(\vct x_g/R\).}

\entrada{nucleo}{Núcleo de influencia \(U\) o \(T\)}{3, 4, 6}{
Es una función matricial que indica cómo se relacionan una fuente
\(\vct\xi\) y un punto \(\vct x\). \(U\) es el desplazamiento del estado
fundamental; \(T\), su tracción. Se evalúan para muchos pares y se
integran para construir \(G\) y \(H\). No son, por sí solos, las
matrices globales \(G\) y \(H\).}

\entrada{numeroonda}{Número de onda \(k\)}{2, 4}{
Relaciona frecuencia y velocidad mediante \(k_s=\omega/c_s\) y
\(k_p=\omega/c_p\), con unidades de inverso de longitud. Puede ser
complejo en un material amortiguado. No es el número de nodos ni el
índice de una componente.}

\grupo{O--R}

\entrada{ondas}{Ondas P y S}{2, 4}{
Son dos modos de propagación elástica. En una onda P el movimiento es
longitudinal; en una S es transversal. Sus velocidades son \(c_p\) y
\(c_s\). Ambas aparecen en la solución fundamental elástica, incluso
si la respuesta global del ejemplo es radial.}

\entrada{operador}{Operador}{1, 3, 4}{
Es una regla matemática que transforma un campo en otro. Por ejemplo,
el operador de tracción transforma el desplazamiento auxiliar en su
tracción sobre una superficie orientada. Los operadores pueden contener
derivadas, integrales o multiplicaciones por coeficientes.}

\entrada{particion}{Partición de la unidad}{5, 6}{
Es la propiedad \(\sum_{a=1}^3P_a(\eta)=1\) dentro de cada elemento.
Garantiza que un campo nodal constante se interpole como constante.
Por eso es importante para deducir la identidad de traslación rígida
utilizada en \(H^s\).}

\entrada{peso}{Peso de cuadratura}{5, 6}{
El peso \(w_q\) indica cuánto aporta el punto de Gauss \(q\) a una
integral de referencia. Para integrar un nodo local se combina con
geometría y forma:
\(W_{qa}=P_a(\eta_q)J_qw_q\). Aunque los pesos ordinarios de Gauss son
positivos, \(W_{qa}\) puede ser negativo porque \(P_a\) puede serlo.}

\entrada{presion}{Presión \(p\)}{1, 7}{
Es una carga escalar normal por unidad de área, en Pa. En el ejemplo
\(p=100\) Pa representa compresión: la tracción vectorial es
\(\vct t=-p\vct n\). Sus componentes cartesianas cambian alrededor
del círculo, aunque la presión sea uniforme.}

\entrada{punto}{Punto de integración \(\vct x_q\)}{3--6}{
Es una posición del contorno donde se evalúan los núcleos durante la
cuadratura. Para una fuente fija \(\vct\xi\) se recorren los elementos
y sus puntos \(\vct x_q\). En esta implementación la cuadratura
adaptada evita evaluar exactamente en la fuente.}

\entrada{reciprocidad}{Reciprocidad de Betti}{3}{
Es una simetría que relaciona dos estados de un mismo material lineal.
Al comparar el estado real y el auxiliar a la misma frecuencia permite
obtener la identidad de contorno. Para el material complejo de esta
formulación se usan productos bilineales, sin conjugación compleja.}

\entrada{regularizacion}{Regularización de la singularidad}{4, 6}{
Es reorganizar una expresión singular para poder tratarla de manera
controlada. Aquí se separa \(T=T^s+(T-T^s)\): la diagonal estática se
obtiene por traslación rígida y la corrección dinámica se integra.
La resta elimina la singularidad fuerte, pero puede perder precisión
numérica muy cerca de la fuente.}

\entrada{residuo}{Residuo algebraico}{7}{
Mide cuánto se incumple el sistema ensamblado después de resolverlo:
\(\vct r_{\mathrm{alg}}=H\vct u-G\vct t\). Un residuo pequeño indica
que se resolvió bien \emph{ese sistema}; no garantiza que malla,
núcleos y cuadratura aproximen correctamente el problema físico.}

\entrada{resonancia}{Resonancia}{7}{
Es una región de frecuencia en la que la respuesta del sistema puede
amplificarse considerablemente. La ubicación y el tamaño de la respuesta
dependen del material, la geometría, la carga y el amortiguamiento.
Cerca de una resonancia conviene revisar malla, cuadratura y
condicionamiento del sistema.}

\grupo{S--V}

\entrada{singularidad}{Singularidad}{3, 4, 6}{
Es un comportamiento no acotado del núcleo cuando el punto de integración
se acerca a la fuente: \(r\to0\). Aquí \(U\) tiene comportamiento
logarítmico integrable y \(T\) presenta una parte del orden de \(1/r\).
La singularidad requiere tratamiento matemático y numérico; no se
resuelve simplemente asignando cero al núcleo en \(r=0\).}

\entrada{sistema}{Sistema lineal \(H\vct u=G\vct t\)}{6, 7}{
Es el conjunto finito de ecuaciones que resulta de discretizar y
ensamblar la identidad de contorno. En el ejemplo, la carga conocida
forma \(\vct b=G\vct t\) y se resuelve \(H\vct u=\vct b\).
No es necesario construir explícitamente la inversa de \(H\).}

\entrada{solucionfundamental}{Solución fundamental}{3, 4}{
Es la respuesta matemática de un medio auxiliar a una fuente puntual
unitaria. En este caso genera el núcleo de desplazamiento \(U\) y,
al aplicar el operador de tracción, el núcleo \(T\).
No tiene que satisfacer por sí sola la presión del cilindro: eso lo
hace la formulación integral con las condiciones reales.}

\entrada{tangente}{Tangente \(\vct\tau\)}{5}{
Es la dirección que sigue la curva al recorrerla. Si
\(\vct x'=(x',y')\), la tangente unitaria es
\(\vct\tau=\vct x'/J\). Para la malla antihoraria usada aquí, la normal
exterior se obtiene rotándola hacia la derecha:
\(\vct n=(y',-x')/J\).}

\entrada{tensor}{Tensor}{1, 3, 4}{
Es un objeto cuyas componentes relacionan direcciones. El esfuerzo
\(\boldsymbol\sigma\) y los núcleos \(U,T\) se representan mediante
matrices de componentes en dos dimensiones. Una matriz de componentes
no es necesariamente la matriz global ensamblada de todo el problema.}

\entrada{terminolibre}{Término libre \(c(\vct\xi)\)}{3, 6}{
Es la contribución local que aparece al llevar la fuente interior al
contorno en la identidad integral. En una frontera suave vale
\(c=\tfrac12 I\); en esquinas depende de la geometría local. La diagonal
estática construida por movimiento rígido ya lo incorpora: no se añade
otro \(\tfrac12 I\).}

\entrada{traccion}{Tracción \(\vct t\)}{1, 3, 7}{
Es el vector de fuerza por unidad de área que actúa sobre una superficie
con normal \(\vct n\): \(\vct t=\boldsymbol\sigma\vct n\), en Pa.
Depende de la orientación de la superficie. En este ejemplo
\(\vct t=-p\vct n\), dirigida hacia el interior.}

\entrada{valorprincipal}{Valor principal de Cauchy, \(\VP\)}{3, 6}{
Es una forma prescrita de tomar el límite de una integral singular,
excluyendo una pequeña vecindad de la fuente y combinando los aportes
antes de hacerla tender a cero. Por ejemplo,
\(\VP\int_{-1}^{1}\dd s/s=0\) con exclusión simétrica.
No equivale a no evaluar el integrando en un solo punto.}

\entrada{vector}{Vector}{1, 6, 7}{
Es una lista ordenada de componentes. En un nodo,
\(\vct u_g=(u_{g,x},u_{g,y})\); en el sistema global se apilan las
componentes de todos los nodos. Un vector tiene dirección y magnitud en
el plano; la lista global también organiza las incógnitas algebraicas.}

\entrada{velocidad}{Velocidad de onda \(c_s,c_p\)}{2, 4}{
Caracteriza las ondas de corte y longitudinales:
\(c_s=\sqrt{\mu^*/\rho}\),
\(c_p=\sqrt{(\lambda^*+2\mu^*)/\rho}\).
En el modelo amortiguado pueden ser complejas. Para interpretar
propagación y atenuación se examina el número de onda complejo,
no solo la parte real de \(c\).}

\newpage
\section*{Ejemplo guiado: seguir una influencia de principio a fin}
\begin{center}
\begin{tikzpicture}[>={Stealth[length=2mm]},scale=0.93]
 \fill[blue!4] (0,0) circle(1.7);
 \draw[thick,blue!55!black] (0,0) circle(1.7);
 \fill[blue!70!black] (1.7,0) circle(2.4pt)
   node[below right] {\(\vct\xi=\vct x_0\)};
 \fill[orange!80!black] (1.20,1.20) circle(2.4pt)
   node[above left] {\(\vct x\)};
 \draw[->,thick] (1.63,0.12)--(1.28,1.08)
   node[midway,left] {\(r\)};
 \draw[->,thick,teal!65!black] (1.20,1.20)--(1.8,1.8)
   node[above right] {\(\vct n(\vct x)\)};
 \draw[->,thick,red!65!black] (0,2.25)--(0,1.75)
   node[pos=0,above] {\(p\)};
 \node at (0,-0.5) {\(\Omega\)};
 \node at (-1.9,-1.25) {\(\Gamma\)};
 \node[anchor=west,align=left,text width=7cm] at (4.1,0.1) {
   \textbf{Fuente:} se fija una ecuación.\\[2mm]
   \textbf{Integración:} se recorre el borde.\\[2mm]
   \textbf{Núcleos:} conectan ambos puntos.\\[2mm]
   \textbf{Presión real:} actúa en todo \(\Gamma\).};
\end{tikzpicture}
\end{center}

Supongamos \(R=6\) m, \(p=100\) Pa y \(n_e=8\):
\begin{enumerate}
 \item Hay \(N=16\) nodos globales y \(2N=32\) componentes escalares.
 Elige el nodo global 0 como fuente: \(\vct\xi=\vct x_0=(6,0)\).
 \item Recorre un elemento. El punto \(\vct x\) del dibujo representa
 una posición genérica de integración; en el cálculo se usan los puntos
 de cuadratura \(\vct x_q\). La distancia \(r=\norm{\vct x-\vct\xi}\)
 cambia con cada punto.
 \item Evalúa \(U(\vct x,\vct\xi)\), \(T(\vct x,\vct\xi)\) y la normal
 local. Multiplica por las funciones de forma, el jacobiano y los pesos
 para obtener bloques de influencia del elemento.
 \item El nodo global 2 pertenece al elemento 0 como nodo local 3 y al
 elemento 1 como nodo local 1. Sus dos contribuciones se suman en el
 mismo bloque global \((0,2)\).
 \item En el nodo 0 la normal circular es \((1,0)\), por lo que la
 tracción prescrita allí es \((-100,0)\) Pa. Se forma la carga de
 \emph{todos} los nodos, se resuelve y se proyecta radialmente.
\end{enumerate}

\begin{resultado}{Un resultado concreto para conectar los términos}
Con \(\rho=100\ \mathrm{kg/m^3}\), \(\mu_0=10^6\) Pa, \(\nu=0.25\),
\(\beta=0.05\) y \(\omega=40\ \mathrm{rad/s}\), la respuesta calculada
en el nodo 0 es aproximadamente
\[
 u_{r,0}=-2.395966\times10^{-4}
          +\ii\,4.085472\times10^{-5}\ \mathrm m.
\]
Su amplitud es \(|u_{r,0}|\simeq0.243055\) mm. La amplitud es positiva;
el signo y el desfase se interpretan mediante la respuesta compleja.
\end{resultado}

\begin{alerta}{Cuatro distinciones que evitan errores}
\textbf{Fuente \(\ne\) carga física.}
\textbf{Local \(\ne\) global.}
\textbf{Núcleo \(\ne\) matriz ensamblada.}
\textbf{Residuo pequeño \(\ne\) error físico pequeño.}
\end{alerta}

\newpage
\section*{Cómo leer la ecuación de contorno}
Para una fuente fija \(\vct\xi\) situada en el borde, se escribe
\[
 \underbrace{c\,\vct u(\vct\xi)}_{\text{término local}}
 +\underbrace{\VP\int_\Gamma T(\vct x,\vct\xi)\vct u(\vct x)\dd\Gamma}
   _{\text{influencia del desplazamiento desconocido}}
 =\underbrace{\int_\Gamma U(\vct x,\vct\xi)\vct t(\vct x)\dd\Gamma}
   _{\text{influencia de la tracción conocida}}.
\]
\begin{enumerate}
 \item \(\Gamma\): se suma la contribución de todo el contorno.
 \item \(\vct\xi\) es fija; \(\vct x\) cambia durante la integración.
 \item \(U\) y \(T\) son funciones de influencia, no las incógnitas.
 \item \(\VP\) prescribe el límite singular; \(c\) es el término libre.
 \item Al interpolar y ensamblar:
 \(\boxed{H(\omega)\vct u=G(\omega)\vct t}\).
\end{enumerate}

\begin{idea}{Convención de índices de los núcleos}
En estas guías, \(U_{ij}\) es la componente \(j\) del desplazamiento
auxiliar producido por una fuerza unitaria en dirección \(i\).
\(T_{ij}\) es la componente \(j\) de la tracción de ese mismo estado:
se aplica el operador de tracción a la fila \(U_{i\cdot}\).
En la integral se suman los productos \(U_{ij}t_j\) y \(T_{ij}u_j\).
Aunque \(U\) sea simétrica, \(T\) no tiene por qué serlo.
\end{idea}

\section*{Símbolos que conviene tener a mano}
\begin{center}
\small
\begin{tabular}{p{0.20\linewidth}p{0.70\linewidth}}
\toprule
Símbolo & Qué significa en este problema\\
\midrule
\(\Omega,\Gamma\) & Región de material y su frontera.\\
\(R,r,\vct e\) & Radio del cilindro; distancia y dirección desde la fuente.\\
\(\vct\xi,\vct x\) & Fuente fija; punto variable de integración.\\
\(n_e,N\) & Número de elementos; número de nodos globales (\(N=2n_e\)).\\
\(e,a,g,q\) & Elemento, nodo local, nodo global y punto de cuadratura.\\
\(\eta,J,w_q,W_{qa}\) & Coordenada local, jacobiano y pesos de integración.\\
\(\vct n,p,\vct t\) & Normal exterior, presión y tracción.\\
\(\rho,\nu,\beta\) & Densidad, coeficiente de Poisson y amortiguamiento.\\
\(\lambda^*,\mu^*\) & Módulos elásticos complejos.\\
\(c_s,c_p,\omega\) & Velocidades S y P; frecuencia angular.\\
\(z_s,z_p,\gamma,\psi,\chi\) & Argumentos, relación de velocidades y
funciones radiales de los núcleos.\\
\(U,T,T^s\) & Núcleos de desplazamiento, tracción dinámica y tracción estática.\\
\(c,C_c\) & Término libre; no son las velocidades \(c_s,c_p\).\\
\(H,G,H^s\) & Matrices que multiplican desplazamientos y tracciones;
referencia estática \(H^s\).\\
\(\vct u,u_r\) & Desplazamiento cartesiano y su proyección radial.\\
\bottomrule
\end{tabular}
\end{center}

\begin{alerta}{Una misma letra puede cambiar de significado}
\(i\) puede ser una componente o la unidad imaginaria \(\ii\);
\(t\) puede ser tiempo, mientras \(\vct t\) es tracción;
\(J\) puede ser jacobiano y \(J_n\) una función de Bessel;
\(\gamma\) puede denotar deformación cortante o la relación \(c_s/c_p\).
El contexto y los subíndices indican qué se está usando.
\end{alerta}

\par\medskip{\footnotesize
\textbf{Documentos relacionados.} Este diccionario complementa los
PDF de los pasos 1--7 de esta carpeta y el esquema
\texttt{../esquema\_ecuaciones\_bem.pdf}. Para recompilar el archivo
\texttt{.tex}, conservar \texttt{estilo.tex} en la misma carpeta y
ejecutar \texttt{pdflatex} dos veces desde allí.}
\end{document}
